\documentclass[11pt,a4paper]{article}
\usepackage[utf8]{inputenc}
\usepackage{amsmath,amssymb,amsthm}
\usepackage[margin=1in]{geometry}

\newtheorem{theorem}{Theorem}
\newtheorem{lemma}[theorem]{Lemma}

\title{Problem 127: abc-hits}
\author{Project Euler Solution}
\date{}

\begin{document}
\maketitle

\section{Problem Statement}

A triple $(a, b, c)$ is an \emph{abc-hit} if $\gcd(a,b) = 1$, $a < b$, $a + b = c$, and $\mathrm{rad}(abc) < c$. Find $\sum c$ for all abc-hits with $c < 120000$.

\section{Mathematical Foundation}

\begin{theorem}[Pairwise coprimality]
If $\gcd(a,b) = 1$ and $a + b = c$, then $a$, $b$, $c$ are pairwise coprime.
\end{theorem}

\begin{proof}
If $p \mid a$ and $p \mid c$, then $p \mid (c - a) = b$, contradicting $\gcd(a,b) = 1$. Similarly $\gcd(b,c) = 1$.
\end{proof}

\begin{theorem}[Radical factorization]
If $a$, $b$, $c$ are pairwise coprime, then $\mathrm{rad}(abc) = \mathrm{rad}(a) \cdot \mathrm{rad}(b) \cdot \mathrm{rad}(c)$.
\end{theorem}

\begin{proof}
Pairwise coprimality means the prime divisors of $a$, $b$, $c$ are disjoint. Hence $\prod_{p \mid abc} p = \prod_{p \mid a} p \cdot \prod_{p \mid b} p \cdot \prod_{p \mid c} p$.
\end{proof}

\begin{lemma}[Reformulated condition]
The abc-hit condition becomes $\mathrm{rad}(a) \cdot \mathrm{rad}(b) < c / \mathrm{rad}(c)$.
\end{lemma}

\begin{proof}
Divide $\mathrm{rad}(a) \cdot \mathrm{rad}(b) \cdot \mathrm{rad}(c) < c$ by $\mathrm{rad}(c) > 0$.
\end{proof}

\begin{lemma}[Coprimality via radicals]
$\gcd(a,b) = 1$ iff $\gcd(\mathrm{rad}(a), \mathrm{rad}(b)) = 1$.
\end{lemma}

\begin{proof}
A prime $p$ divides $\gcd(a,b)$ iff it divides both $\mathrm{rad}(a)$ and $\mathrm{rad}(b)$.
\end{proof}

\begin{theorem}[Early termination]
Iterating $a$ in order of increasing $\mathrm{rad}(a)$, once $\mathrm{rad}(a) \ge c/\mathrm{rad}(c)$, no further $a$ can yield an abc-hit for that $c$.
\end{theorem}

\begin{proof}
$\mathrm{rad}(a) \cdot \mathrm{rad}(b) \ge \mathrm{rad}(a) \ge c/\mathrm{rad}(c)$, violating the condition.
\end{proof}

\section{Algorithm}

\begin{enumerate}
\item Sieve $\mathrm{rad}(n)$ for $n < N = 120000$.
\item Sort indices by radical.
\item For each $c = 3, \ldots, N-1$: iterate $a$ by increasing radical, set $b = c - a$, check $a < b$, $\gcd(\mathrm{rad}(a), \mathrm{rad}(b)) = 1$, and $\mathrm{rad}(a) \cdot \mathrm{rad}(b) < c/\mathrm{rad}(c)$.
\item Sum qualifying $c$ values.
\end{enumerate}

\section{Complexity Analysis}

\begin{itemize}
\item \textbf{Time:} $O(N \log \log N)$ for sieve, $O(N \log N)$ for sort and main search (empirical).
\item \textbf{Space:} $O(N)$.
\end{itemize}

\section{Answer}

\[
\boxed{18407904}
\]

\end{document}
